\begin{statement}{E.25}{Proved}{The complete joint quartic distribution test}
Keep $B,F,N_B,G_6$ and the twenty-four retained words $f_i\in F'$
of E.1, indexed $i=0,\ldots,23$ in the saved certificate order.
Let $q_{24}:F\to\mathrm{SL}_2(\mathbb Z/24)$ send
\[
 x\longmapsto U^2=\begin{pmatrix}1&2\\0&1\end{pmatrix},
 \qquad y\longmapsto W^2=\begin{pmatrix}1&0\\-2&1\end{pmatrix}.
 \tag{E.25.1}
\]
Put $\mathcal N=\{2,3,4,6,12\}$ and retain every coefficient
$c_N$ of D.5 modulo twelve.  For each $i$, the complete set of
pairs $(c_2,c_3)\in(\mathbb Z/12)^2$ attained by words
$f_i k$, $k\in N_B'$, satisfying simultaneously
\[
\begin{gathered}
 q_{24}\delta_N(f_i k)=q_{24}(y)^{c_N(f_i k)}q_{24}(f_i k)
       \quad(N\in\mathcal N),\\
 c_4=2c_2,\qquad c_6=c_2+c_3,\qquad c_{12}=2c_2+c_3
       \pmod {12}
\end{gathered}
 \tag{E.25.2}
\]
is exactly $\{a_i\}\times\mathbb Z/12$, where
\[
\begin{aligned}
 (a_0,\ldots,a_{11})&=(0,6,0,6,6,0,6,0,4,10,4,10),\\
 (a_{12},\ldots,a_{23})&=(10,4,10,8,4,10,4,10,4,10,4,10).
\end{aligned}
 \tag{E.25.3}
\]
For the auxiliary image allowing every correction in
$K_6=N_B'N_B^6$, the exact set is
$\{a_i,a_i+6\}\times\mathbb Z/12$.
The latter calculation does not impose an additional $F'$ condition
on that larger correction group; the former already supplies actual
words in $F'$ with unchanged $G_6$ image.

The joint correction images of $N_B'$ and $K_6$ have orders
\[
             2\,717\,908\,992,
       \qquad65\,229\,815\,808,
 \tag{E.25.4}
\]
respectively.  Thus all twelve retained CM-compatible points and
all twelve retained extra points pass this joint test.  Each admits
every value of $\rho_3=-c_3\pmod4$; restricting both $c_2,c_3$ to
even residues still leaves both values $0,2$ of $\rho_3\pmod4$.
This assertion concerns exactly the finite identities (E.25.2).
It asserts neither generation in their induced cover quotient nor
a full profinite GT realization.
\end{statement}
\begin{proof}
The mod-four images of $U^2,W^2$ generate the four-group $V_4$
of E.1.  Choosing the unique global sign with reduction in that
four-group identifies the image in (E.25.1) with
\[
 H_{24}=G_2\times\mathrm{SL}_2(\mathbb F_3),\qquad
 G_2=\{g\in\mathrm{SL}_2(\mathbb Z/8):g\bmod4\in V_4\},
 \qquad |G_2|=32.
 \tag{E.25.5}
\]
Reduction projectively modulo twelve is the original map on $F$.
Consequently $q_{24}(N_B)$ lies in the elementary abelian kernel
of reduction from modulus twenty-four to modulus twelve.
It follows that $q_{24}$ kills both $N_B'$ and $N_B^6$.
Thus the last matrix in (E.25.2) is the fixed matrix $q_{24}(f_i)$.

Write $H_N=\ker(F\to\mathbb Z/N)$ as in D.5 and
$A=H_{24}\times\mathbb Z/12$.  The actual homomorphism
\[
 \Psi_N:H_N\longrightarrow A,\qquad
 w\longmapsto(q_{24}\delta_N(w),c_N(w)\bmod12)
 \tag{E.25.6}
\]
is induced to $F$ using the transversal $1,x,\ldots,x^{N-1}$.
For a word $w$ with $x$ exponent $s$ modulo $N$, its $j$th
coordinate is the image under (E.25.6) of
$x^jwx^{-(j+s\bmod N)}$.  Multiplication shifts the coordinates
by $s$.  This defines an actual homomorphism to the finite wreath
product, with the action convention fixed in the program.

Here is the exact reduction used to compute its joint derived
image.  The image of the $x$ exponent of $N_B$ in $\mathbb Z/12$
is the even subgroup of order six.  Put $N_0=N_B\cap H_{12}$.
Choose actual words $t_s\in N_B$ for the six shifts.  If
$n_1,\ldots,n_{97}$ is the literal basis from E.1, the elements
\[
      t_s n_j t_{s+\operatorname{exp}_x(n_j)}^{-1}
 \tag{E.25.7}
\]
generate $N_0$.  On $N_0$, taking coordinate zero in each of
the five induced representations is a homomorphism into $A^5$.
Furthermore $N_B'$ is the normal closure in $N_0$ of
\[
             t_s[n_i,n_j]t_s^{-1}
       \quad(s\text{ among the six shifts},\ i<j).
 \tag{E.25.8}
\]
Indeed the ordinary normal closure of the basis commutators in
$N_B$ is $N_B'$, and every conjugating element is a product of
an element of $N_0$ and one of the $t_s$.  This proves the same
assertion after taking closures and then any finite image.
Thus the normal closure of the images of (E.25.8) under the
image of (E.25.7) is exactly the desired joint correction group.
It retains all dependencies between the five coordinates.
Adjoining the images of $n_j^6$ gives the image of $K_6$:
modulo $N_B'$ the source is abelian, so these powers generate
all sixth powers there.  Their shifts are zero since every
$n_j$ has even $x$ exponent.

For a faithful finite model, write each member of $G_2$ uniquely
as $\bar x^a\bar y^b[\bar x,\bar y]^e$, with
$a,b\in\mathbb Z/4$ and $e\in\mathbb Z/2$.
Its two cyclic-four quotient actions, together with its action
on the four cosets of $\langle\bar x,\bar y^2\rangle$, give a
faithful permutation representation of degree twelve.
The other factors act on the eight nonzero vectors of
$\mathbb F_3^2$ and on disjoint cyclic sets of sizes four and
three.  Hence $A$ has the explicit faithful degree-twenty-seven
representation used in the program.  Its five induced actions
have total degree $27(2+3+4+6+12)=729$; the shift-zero image
has degree $135$.

The program \texttt{bc\_distribution\_quartic.py} verifies every
multiplication in the order-thirty-two model.  The 385 projected
Schreier generators reduce, by exact membership, to eleven
generators.  The 11,419 distinct projected commutator seeds
reduce to seventeen; closure under all eleven conjugating
generators is checked before its order is computed.  Every
generator retained or added is the image of the displayed
actual source words.  Exact Schreier--Sims membership gives
(E.25.4).

For each $i$, the program enumerates all 144 pairs $(a,b)$,
substitutes $(c_2,c_3,c_4,c_6,c_{12})=(a,b,2a,a+b,2a+b)$,
and tests the resulting single correction tuple for membership
in each of the two joint groups.  Its complete records give
(E.25.3) and the auxiliary sets stated above.  A tuple in the
first group has one preimage $k\in N_B'$, so it supplies one
word satisfying all five identities.  The three-cover instance
was also computed in the full induced ambient group before
the shift-zero reduction was introduced; its complete records
agree, with correction orders 192 and 576.  Source, output,
literal input words and generator arrays are saved in
\texttt{code/bc\_distribution\_quartic*} and
\texttt{code/bc\_cm\_literal\_words.json}.

The convention $\rho_N=-c_N$ is the one proved in D.5.
No identification of this coefficient on the whole GT source
with a different matrix-defined extension of a Galois Kummer
character is needed for this finite calculation.
\end{proof}

\begin{statement}{E.26}{Proved}{The exact Fermat proper-Tate quotient and its boundary refinement}
Start with the discrete group $B=C_2*C_3$.  Use the literal
level-two generators of the E.1 verifier,
$x=(r^{-1}s)^2=(sr)^{-2}$ and $y=(sr^{-1})^2=sxs$.
Let $F=\langle x,y\rangle\triangleleft B$, and put
\[
 H=\ker(F\to(\mathbb Z/3)^2),\qquad
 T=H^{\rm ab}/\langle\text{the nine cusp classes}\rangle.
 \tag{E.26.1}
\]
The based cubic cover identifies $H$ with the topological
fundamental group of the punctured Fermat cubic, and
$T\simeq\mathbb Z^2$ is its proper integral homology lattice.
After profinite completion the corresponding finite cover
identifies $\widehat H$ with the geometric \'etale fundamental
group, and $\widehat T=\varprojlim_m T/mT$
with the full Tate module of its smooth compactification.
For $m\ge1$, set $J_m=\ker(H\to T/mT)$; it is normal in $B$.
There are actual quotient maps
\[
       B\longrightarrow G_{3m}\longrightarrow B/J_m.
 \tag{E.26.2}
\]
In particular every proper two-power Tate quotient of this
specified cover is dominated by the homology tower of E.4.
The inverse-limit comparison is the surjection
\[
       \widehat B/\overline{[N_B,N_B]}
            \longrightarrow\varprojlim_m B/J_m.
 \tag{E.26.3}
\]
These statements retain the intrinsic Fermat curve and its
based maps; no CM twist comparison is implicit in them.

The full rank-ten homology quotient
$B/[H,H]H^4$ gives additional information.  In the ordered
basis $\eta_1,\ldots,\eta_{10}$ specified below, the first two
literal $N_B$ generators satisfy
\[
 [n_1,n_2]\longmapsto2\eta_3-2\eta_{10}\ne0
            \quad\hbox{in }H^{\rm ab}/4,
 \tag{E.26.4}
\]
where $n_1=y^6$ and $n_2=yx^{-1}yx^{-2}yx^{-1}y$.
Consequently the actual common refinement
\[
     B/(K_6\cap[H,H]H^4)\longrightarrow G_6
 \tag{E.26.5}
\]
cannot be dominated by any $G_j$ through the specified maps
from $B$.
\end{statement}
\begin{proof}
Use the transversal $x^a y^b$, $0\le a,b<3$, and order its
nontrivial Schreier edges first by $(a,b)$ and then by $x,y$.
The cover is obtained by adjoining $t^{1/3}$ and
$(1-t)^{1/3}$; its compactification is $u^3+v^3=1$.
At the standard tangent choose the former Puiseux branch and
the latter branch with constant coefficient one.
The finite-index subgroup $H$ has the full induced profinite
topology: the core in $B$ of any finite-index subgroup of $H$
still has finite index.  The finite-cover comparison therefore
gives precisely the completion and filling maps in the statement.
The ten nontrivial Schreier edges are the generators used in
\texttt{de\_fermat\_schreier.py}.  Integral rewriting gives
the primitive quotient matrix
\[
 T_0=\begin{pmatrix}
 -1&-1&0&0&1&0&0&1&0&0\\
  1&0&0&-1&-1&0&0&0&1&0
 \end{pmatrix}.
 \tag{E.26.6}
\]
The nine cusp columns have Smith factors $1^8,0$, and $T_0$
has Smith factors $1,1$.  Thus (E.26.6) is the actual integral
filling map with no finite-index correction.  Conjugation by
either $x$ or $y$ induces on this quotient the same matrix
\[
          R=\begin{pmatrix}-1&-1\\1&0\end{pmatrix},
           \qquad R^2+R+1=0.
 \tag{E.26.7}
\]
These identities are checked by rewriting each of the ten
conjugated basis words over $\mathbb Z$.  They hold after
every finite coefficient reduction and hence on $\widehat T$.
Conjugation by $s$ interchanges $x,y$, while conjugation by
$r$ sends $x\mapsto y$, $y\mapsto(xy)^{-1}$.
For precision, the literal identities follow from
\[
 r(r^{-1}sr^{-1}s)r^{-1}=sr^{-1}sr^{-1},\qquad
 r(sr^{-1}sr^{-1})r^{-1}=rsr^{-1}sr=(xy)^{-1}.
\]
For the simultaneously inverted generators
$x_+=x^{-1}=(sr)^2$, $y_+=y^{-1}$, the second formula is
$ry_+r^{-1}=x_+^{-1}(x_+y_+)^{-1}x_+$.
These are the actual changes of convention.
The two conjugations preserve the cubic-cover subgroup and permute its cusp
classes.  This proves the asserted normality of $J_m$.

The reduction map
$Q\to\mathrm{PSL}_2(\mathbb F_3)\simeq A_4\to A_4/V_4$
is a cyclic-three character.  Its restrictions to $x$ and $y$
coincide and are nontrivial: the two matrices are conjugate
order-three transvections, whose projective classes lie in
the same nontrivial coset.  Hence, for every $n\in N_B$,
\[
        \operatorname{exp}_x(n)+\operatorname{exp}_y(n)=0
                    \pmod3.
 \tag{E.26.8}
\]
By (E.26.7), $N_B$ acts trivially on $T$.  Its image in
$F/H$ lies in the cyclic line $\{(a,-a):a\in\mathbb Z/3\}$.
The image of $N_B\cap H$ in $B/J_m$ is therefore central
in the image of $N_B$, and the quotient by it is cyclic.
A group with a central subgroup and cyclic quotient is
abelian: one lift of a cyclic generator and the central
subgroup generate it.  It follows that $N_B'\subset J_m$.
Moreover $n^3\in H$ for $n\in N_B$, so $n^{3m}\in J_m$.
These two inclusions prove exactly (E.26.2).  All its targets
are finite, so the same maps extend uniquely to $\widehat B$.
Their kernels contain the closure of $[N_B,N_B]$, and taking
the inverse limit over divisibility of $m$ proves (E.26.3).
For the simultaneous comparison
with $G_6$, use $G_{\operatorname{lcm}(6,3m)}$.

For the last assertion, both $n_1,n_2$ are literal basis
members of the kernel of $B\to Q$.  Rewriting their
commutator in the ten-generator cubic-cover basis gives
the integral vector $(0,0,2,0,0,0,0,0,0,-2)$.
Its image under $T_0$ is zero but its reduction modulo four
is nonzero.  Thus $N_B'$ is not contained in $[H,H]H^4$,
nor in the kernel in (E.26.5), whereas every $G_j$ kills
$N_B'$.  A factorization through any specified $G_j$
would kill this particular word, proving the claimed
failure of factorization.
The program \texttt{bc\_cm\_fermat\_factorization.py}
retains the integral deck matrices and all 4,656 basis
commutator images.  Each has zero proper image, and the
displayed boundary witness is its first nonzero reduction
modulo four.  It also checks that translation by $(1,-1)$
has $\operatorname{rank}(A-1)=6$ on the full rank-ten module.
Its output, JSON and exact arrays are saved alongside it.
\end{proof}
